\section{Análisis}

Firmware \texttt{dsPicDaq} con el PLL ($F_{cy} = 50$\,MHz) y reporte de la
causa de reinicio en los cuatro bits altos del status. Comparado con el
firmware a 4\,MHz (\texttt{resultados/2026-10-05/dspic\_1khz}):

\begin{table}[H]
\centering
\begin{tabular}{@{}lrr@{}}
\toprule
 & $F_{cy} = 4$\,MHz & $F_{cy} = 50$\,MHz \\
\midrule
CRC incorrecto en la PC        & 83.7\,\% & 0.68\,\% \\
Respuestas de arranque         & 5.9\,\%  & 0 \\
Reinicios reportados           & --       & 0 \\
Longitud incorrecta (dsPIC)    & 0.07\,\% & 0.72\,\% \\
\bottomrule
\end{tabular}
\end{table}

El bit 7 de los bytes ya llega correcto: con $F_P = 50$\,MHz el SPI esclavo
presenta el primer bit de cada byte mucho antes del flanco en que lo muestrea
el FT2232H. Las respuestas de arranque (secuencia 0) desaparecieron y el
dsPIC no reportó ningún reinicio, por lo que las del firmware anterior
estaban ligadas al reloj de 4\,MHz; con el reporte de causa, un reinicio
futuro ya no pasaría inadvertido. La única respuesta contada como de
arranque es una trama legítima con eco de la secuencia 0.

Queda un defecto en el 0.7\,\% de las transferencias. En 672 de las 675
respuestas con CRC incorrecto, el dsPIC reporta en la respuesta siguiente
que esa misma transferencia no tuvo 8 bytes: la falla daña ambos sentidos a
la vez. Comparando con la respuesta esperada, la corrupción empieza en el
byte 1 en el 63\,\% de los casos y en el byte 7 en el 26\,\%.

El byte 1 es el que \texttt{armar\_dma()} entrega con una solicitud por
software cuando el evento SPITBE ocurrió antes de habilitar el canal; el
README ya lo señalaba como el punto que no podía verificarse sin hardware.
Con la CPU 12.5 veces más rápida cambió la relación de tiempos entre esa
solicitud y el evento, y en algunas tramas el canal de transmisión parece
recibir un disparo de más o de menos. Los errores en el byte 7 apuntan a la
frontera con el flanco de subida de \texttt{CS}. Ambos casos deben
observarse con el analizador lógico (SCK, MISO, \texttt{CS} y RD10) antes de
modificar \texttt{armar\_dma()}.
